% RESUMO (NBR 6028)
% --------------------------------------------------------------------------
\setlength{\absparsep}{18pt}
\begin{resumo}
A democratização do acesso à baixa órbita terrestre por meio de nanossatélites da classe CubeSat impõe desafios rigorosos de engenharia de sistemas. Durante o lançamento e a ejeção, o chassi estrutural é submetido a carregamentos dinâmicos severos, como o choque transiente na interface do \textit{Poly-Picosatellite Orbital Deployer} (P-POD) e a vibração aleatória acústica regulamentada pela norma NASA GSFC-STD-7000A (GEVS). Chassis monolíticos convencionais apresentam elevada massa e transferem altos níveis de aceleração à carga útil devido ao acoplamento dinâmico direto. Este trabalho propõe uma rota computacional de alto desempenho para a concepção, a modelagem paramétrica e a otimização multiobjetivo de um chassi CubeSat 1U com painéis constituídos por metamateriais auxéticos reentrantes em liga aeroespacial AlSi10Mg, fabricado por Fusão em Leito de Pó a Laser (L-PBF). Um modelo paramétrico tridimensional, sujeito a critérios estritos de Manufatura Aditiva (DfAM), foi integrado a rotinas automatizadas no Ansys MAPDL para gerar um conjunto de 250 simulações de alta fidelidade (análise modal por bloco de Lanczos e resposta espectral à vibração aleatória). Treinou-se um modelo substituto neural residual guiado por física (\textit{Physics-Guided ResNet}), com função de perda composta pelas leis fenomenológicas de Gibson-Ashby e por um termo de monotonicidade de rigidez, que alcançou $\overline{R^2} = \num{0.9874}$ e latência de inferência de \SI{0.405}{\milli\second} (aceleração superior a \num{200000} vezes). Uma máquina de estados finitos determinística, implementada em LangGraph como grafo de fluxo de dados fortemente tipado, orquestrou o descarte precoce de candidatos inviáveis. A otimização multiobjetivo pelo algoritmo NSGA-II, com \num{10000} avaliações de função, mapeou a fronteira de Pareto entre massa e transmissibilidade. Pelo método multicritério TOPSIS, elegeu-se o projeto ótimo de voo ($\theta = \SI{65.75}{\degree}$, $t = \SI{0.613}{\milli\metre}$, $\nu_{\text{eff}} = \num{-13.81}$). A validação numérica de alta ordem (\textit{FEA Ground Truth}) confirmou erro de \SI{0.99}{\percent} na tensão de pico e de \SI{0.008}{\percent} na transmissibilidade. Em relação ao chassi monolítico convencional, a estrutura auxética otimizada proporcionou redução de massa de \SI{62.49}{\percent} (\SI{0.308}{\kilo\gram} $\to$ \SI{0.116}{\kilo\gram}), atenuação elastodinâmica de \SI{76.50}{\percent} na transmissibilidade vibratória (\SI{12.6}{\decibel} por filtragem mecânica e desacoplamento de impedância acústica), frequência fundamental de \SI{579.62}{\hertz} e pleno atendimento aos critérios de escoamento e de fadiga estocástica de Steinberg ($D = \num{0.046} \ll \num{0.25}$), qualificando a tecnologia para operação aeroespacial.

\vspace{\onelineskip}
\noindent
\textbf{Palavras-chave}: CubeSat 1U. Metamateriais auxéticos. \textit{Design for Additive Manufacturing}. Redes neurais guiadas por física. Otimização multiobjetivo NSGA-II. NASA GEVS.
\end{resumo}

